\documentclass{article}
\usepackage[margin=1in]{geometry}
\usepackage[T1]{fontenc}

\begin{document}

\begin{table}[ht]
\centering
\caption{Aggregate compression results excluding 7z.}
\label{tab:compression-summary}
\begin{tabular}{llrrrr}
\hline
Dataset & Technique & Original (bytes) & Compressed (bytes) & Ratio & Time (ms) \\
\hline
Random & OpenZL     & 22,369,280 & 22,369,536 & 1.000011 & 276.915 \\
Random & Gzip       & 22,369,280 & 22,376,165 & 1.000308 & 630.040 \\
Random & Zstandard  & 22,369,280 & 22,369,868 & 1.000026 & 19.565 \\
Random & ZIP        & 22,369,280 & 22,374,766 & 1.000245 & 1185.037 \\
\hline
Trend & OpenZL       & 22,369,280 & 1,454,615 & 0.065027 & 256.617 \\
Trend & Gzip         & 22,369,280 & 5,114,980 & 0.228661 & 946.561 \\
Trend & Zstandard    & 22,369,280 & 3,381,713 & 0.151177 & 91.088 \\
Trend & ZIP          & 22,369,280 & 5,116,907 & 0.228747 & 1263.900 \\
\hline
\end{tabular}
\end{table}

\section{Compression Choice for a Streaming Platform}

The compression method for a streaming platform should be selected according to
the data type and the latency target rather than by compressed size alone.  The
random-data results show that none of the tested methods can meaningfully reduce
unstructured data: each output is approximately the same size as the input.
For this type of payload, compression adds processing and framing overhead, so
the platform should either disable compression for small messages or use a fast
codec only when bandwidth is more important than CPU time.

For structured sensor or telemetry data, the trend-data results are very
different.  OpenZL produced the smallest aggregate output, approximately
6.5\% of the original size, while Zstandard produced approximately 15.1\%.
Zstandard was substantially faster in the tests, with an aggregate time of
91.088~ms compared with 256.617~ms for OpenZL.  This makes Zstandard the best
general-purpose choice for a streaming service where predictable latency,
moderate CPU usage, and broad deployment support are important.  It also has
good support across common programming languages and message brokers.

OpenZL is preferable when the stream has a stable, structured layout and the
main objective is reducing storage or network volume.  Its stronger reduction
on the trend data indicates that it can exploit relationships between values,
but the additional processing cost should be measured against the service's
latency budget.  Compression should normally be performed in batches or
records large enough to amortize codec overhead.  The recommended policy is
therefore to use Zstandard as the default codec, select OpenZL for validated
structured streams with high compression value, and avoid compression for
already-compressed or effectively random payloads.

\section{Kafka versus RabbitMQ for a Priority-Based System}

Kafka and RabbitMQ solve related but different messaging problems.  Kafka is a
distributed append-only event log designed for high-throughput streams,
durable retention, and replay.  RabbitMQ is a broker built around exchanges,
queues, routing rules, acknowledgements, and consumer delivery.  The choice for
a priority-based system should therefore be based on whether priority is a
property of a retained event stream or a property of the next message that a
worker must receive.

\begin{table}[ht]
\centering
\caption{Kafka and RabbitMQ characteristics for priority-based workloads.}
\label{tab:kafka-rabbitmq}
\begin{tabular}{p{0.22\linewidth}p{0.34\linewidth}p{0.34\linewidth}}
\hline
Criterion & Kafka & RabbitMQ \\
\hline
Primary model & Durable, partitioned event log & Routed work queues \\
Priority & Usually implemented through topics, partitions, or separate consumers & Native queue priorities and routing patterns \\
Ordering & Strong ordering within each partition & Ordering within a queue, subject to competing consumers and priority \\
Replay & Built in through retained offsets & Normally requires explicit requeueing or a separate durable design \\
Throughput & Very high for sequential stream processing & Strong, but optimized for brokered delivery and routing \\
Best fit & Analytics, event sourcing, telemetry, and fan-out & Commands, jobs, alerts, and urgent work items \\
\hline
\end{tabular}
\end{table}

Kafka can support priority, but priority is not its fundamental delivery
semantic.  A common design is to publish urgent and normal events to separate
topics or partitions and give urgent consumers more capacity.  Another option
is to maintain separate consumer groups or to use an application-level
scheduler.  These approaches scale well and preserve Kafka's retention and
replay model, but they require extra design work.  A single Kafka partition
also provides ordering, so strict priority can conflict with throughput when a
lower-priority record appears ahead of an urgent record in the same partition.

RabbitMQ is generally the more direct choice when the system must deliver the
highest-priority available job first.  Priority queues, acknowledgements,
dead-letter exchanges, retry policies, and routing keys provide the controls
needed for operational workflows.  For example, an emergency alert can be
routed to a high-priority queue while routine telemetry is handled by a normal
queue.  Consumers can acknowledge work only after processing, allowing failed
messages to be retried without losing them.  The trade-off is that RabbitMQ
does not provide Kafka's equivalent of a long-lived, easily replayable event
history, and very large retained streams are usually better handled by Kafka.

For a priority-based service, RabbitMQ should be selected when the key
requirement is immediate, priority-aware task delivery with per-message
acknowledgement and retry handling.  Kafka should be selected when the system
must retain a large stream, support multiple independent consumers, replay
historical events, or feed analytics and storage pipelines.  A hybrid design
is also appropriate: Kafka can retain the authoritative telemetry or event
stream, while RabbitMQ can carry derived high-priority commands and alerts to
workers.  This separates durable event history from time-sensitive operational
delivery instead of forcing one broker to provide both semantics.

\end{document}